% generated by R/04_tables.R from table4_lineage_groups.md -- do not edit
% needs \usepackage{booktabs,adjustbox} in the preamble
\begin{table}[htbp]
\centering
\caption{Biofilm formation by lineage group}%
\label{tab:biofilmlineage}
\begin{adjustbox}{max width=\linewidth}
\begin{tabular}{lrrl}
\toprule
\textbf{Group} & \textbf{Strains} & \textbf{Wells} & \textbf{Mean OD550} \\
\midrule
sinR & 5 & 40 & 0.042 $\pm$ 0.005 \\
spore & 2 & 16 & 0.144 $\pm$ 0.038 \\
ywcC/epsA-slrR & 3 & 24 & 0.768 $\pm$ 0.118 \\
\bottomrule
\end{tabular}
\end{adjustbox}
\par\smallskip
{\footnotesize\raggedright Blank-corrected OD550 from \texttt{data/biofil.csv}. Groups follow the split an early draft describes: clones carrying a \textit{sinR} mutation never carry one in \textit{ywcC} (BSU\_38220) or upstream of \textit{slrR}, so the likely-vegetative clones divide in two. Means $\pm$ SEM are over wells and are descriptive. Tests treat the strain as the unit of replication -- a linear mixed model of log OD with strain as a random effect, Kenward--Roger degrees of freedom, Tukey-adjusted: sinR vs spore P = 0.769; sinR vs (ywcC/epsA-slrR) P = 0.024; spore vs (ywcC/epsA-slrR) P = 0.151. The \textit{sinR}-versus-spore comparison is the weak one: the spore group is two strains, and they differ 16-fold from each other.\par}
\end{table}
